\documentclass[11pt]{article}
\usepackage{amsmath,amssymb,amsthm,hyperref}
\newtheorem{theorem}{Theorem}
\newtheorem{lemma}[theorem]{Lemma}
\newcommand{\V}{\mathcal V}
\newcommand{\W}{\mathcal W}
\newcommand{\Aug}{\operatorname{AugFam}}
\newcommand{\Fam}{\operatorname{Fam}_{\mathrm{fin}}}
\title{Joint limits and the selected-law augmentation boundary}
\date{16 November 2026}
\begin{document}\maketitle
\paragraph{Status.} The exact text has passed two independent written reviews; an admissibility-class closure qualification found in review is retained explicitly. The target is the next structural gap after the conservative coefficient envelope: whether its new coefficients actually supply native joint context limits. A sufficient construction and an envelope-wide obstruction are distinguished. No general feedback or full extension-contract completion is claimed.

\section{The distinction coefficient completeness does not remove}
The previous manuscript constructs $\W=[\V,\mathbf{Set}]^{op}$, with two opposite Day tensors, from a small original symmetric rig $\V$. Co-Yoneda preserves the original augmented context category and its independent occurrence interaction fully faithfully. The coefficient category $\W$ is complete. This does not automatically imply that its native augmented-family category is complete: native maps aggregate a whole fibre of coefficient objects, rather than supplying an individual arrow from each source coefficient.

There is a second issue. In an exact probabilistic source, a global augmentation still chooses one normalization decomposition. Richer coefficient objects can retain more data without removing that chosen-law invariant. The obstruction below applies to every object of this envelope, not just to an attempted product built on the usual pair carrier.

\section{A general obstruction carried by a selected-law shadow}
Let $\mathbf{FinProb}$ denote finite raw sets with nonnegative weights summing to one and actual functions preserving pushforward exactly. Zero-weight points are retained. For arrows $p:X\to A$ and $q:X\to B$, their induced joint law is $(p,q)_*\mu_X$ on $A\times B$.

\begin{theorem}[Selected-law shadow obstruction]
Suppose a category $\mathcal E$ has functors
\[
 J:\mathbf{FinProb}\longrightarrow\mathcal E,
 \qquad S:\mathcal E\longrightarrow\mathbf{FinProb},
 \qquad SJ\cong\mathrm{id}.
\]
If $A,B$ admit two different joint probability laws with these marginals, then $JA,JB$ have no weak product in $\mathcal E$. In particular they have no categorical product. No assumption that $S$ preserves limits is used, and full faithfulness of $J$ is unnecessary.
\end{theorem}
\begin{proof}
Use the displayed isomorphism to identify $SJ$ with the original objects and maps. A proposed weak product $P$ with projections $a:P\to JA$, $b:P\to JB$ has one joint law
\[
 t=(S(a),S(b))_*\mu_{S(P)}.
\]
For either chosen coupling $\mu$ on $A\times B$, the object $J(A\times B,\mu)$ has the two lifted marginal maps and hence admits a factorization $h$ through $P$. Applying $S$ yields a pushforward-preserving function, and its two projection equations give
\[
 t=(S(a),S(b))_*S(h)_*\mu=(\pi_A,\pi_B)_*\mu=\mu.
\]
Two distinct couplings force two distinct values of the same $t$, a contradiction. Uniqueness of $h$ was not needed.
\end{proof}
For two fair-bit objects, diagonal and antidiagonal fair laws are distinct such couplings. The theorem does not prohibit retaining a category or family of joint apices, one for each compatible choice. It rules out a single universal apex while every new object still has one selected exact law respected by every comparison.

\section{The probabilistic co-presheaf envelope retains this shadow}
Take the original coefficient rig to be the discrete category on the commutative semiring $\mathbb R_{\ge0}$. A covariant presheaf is simply a mass-indexed family of sets $F(s)$. There are no nonidentity base arrows, so the additive Day convolution is
\[
 (F_1\oplus_D\cdots\oplus_DF_n)(s)
 =\coprod_{s_1+\cdots+s_n=s}\prod_iF_i(s_i).
\]
Its additive unit is the singleton presheaf at zero, and the multiplicative unit is the singleton presheaf at one.

An object of $\Aug(\W)$, where $\W=[\mathbb R_{\ge0},\mathbf{Set}]^{op}$, has a finite raw carrier $X$, coefficient presheaves $F_x$, and an augmentation $\bigoplus_xF_x\to y(1)$ in $\W$. Equivalently this is a natural transformation $y(1)\to\bigoplus_DF_x$ in the presheaf category. It chooses an actual tuple
\[
 (p_x,\xi_x)_{x\in X},\qquad \sum_xp_x=1,\quad \xi_x\in F_x(p_x).
\]
This is a selected tuple, not the set of all possible such tuples.

A native map with carrier $f:X\to X'$ has, in the presheaf direction, actual maps
\[
 F'_{x'}\longrightarrow\bigoplus_{fx=x'}F_x.
\]
The augmentation equation says that the chosen target elements are mapped to precisely the chosen source decomposition. Because these maps respect the mass index, it follows that
\[
 p'_{x'}=\sum_{fx=x'}p_x.
\]
Consequently forgetting the selected elements but retaining their indices defines a functor
\[
 S:\Aug(\W)\longrightarrow\mathbf{FinProb}.
\]
On the original scalar representables it is exactly the original raw law. Co-Yoneda therefore supplies the functor $J$ required by the preceding theorem.

\begin{theorem}[The envelope still has no fair-bit weak product]
The two embedded original fair-bit contexts have no weak product in the entire native augmented co-presheaf envelope. Coefficient completeness, nonrepresentable coefficient objects, and extra actual coefficient maps do not remove this obstruction.
\end{theorem}
\begin{proof}
Apply the selected-law shadow theorem to the functor just constructed. The proof does not require guessing the raw carrier or coefficients of a prospective apex.
\end{proof}
The obstruction uses exact selected-law preservation. A law-specification object admitting several complete laws, a retained cone category, or a different comparison doctrine can change that interface. Such a change must be explicit; coefficient enrichment alone has not made it.

\section{When the same envelope does supply native context limits}
There is a positive sufficient condition formulated on the original additive structure rather than on an assumed mixed interaction.

\begin{theorem}[Cocartesian-additive source]
Suppose the original small symmetric rig $\V$ has $\oplus$ as its categorical finite coproduct and $0$ as initial object. Its co-presheaf envelope $\W$ has categorical coproduct as its additive convolution. Its augmented category has finite limits, the raw carrier functor preserves them, and the original augmented source embeds fully faithfully and strongly for its dependent and occurrence structures.
\end{theorem}
\begin{proof}
The additive monoidal structure on $\V^{op}$ is Cartesian. Day convolution of presheaves over a Cartesian monoidal base is pointwise Cartesian product: the hom into a product splits as a product of homs, and the two coends then reduce by Yoneda. Thus on $[\V,\mathbf{Set}]$ the additive convolution is pointwise product, with terminal presheaf as unit; on $\W$ its opposite is categorical coproduct, with initial object as unit.

When the additive structure is a categorical coproduct, an augmentation $\bigoplus_xV_x\to I$ is equivalently one arrow $V_x\to I$ for each $x$. Similarly a fibre aggregation $\bigoplus_{fx=y}V_x\to V'_y$ is equivalently the collection of individual arrows $V_x\to V'_{fx}$. The augmentation equation is precisely their equation over $I$. Hence there is a canonical equivalence
\[
 \Aug(\W)\simeq\Fam(\W/I).
\]
This equivalence retains the entire native arrow, not only its carrier. For the particular envelope, the contraction maps are the canonical coproduct-to-tensor maps: before taking opposites the additive Day product is pointwise product, and the expansion has components obtained by tensoring its projections. After taking opposites, restriction to a coproduct summand therefore gives the expected tensor of coefficient maps. Thus this equivalence also identifies the given dependent and occurrence coefficient operations, rather than silently replacing a possibly different directed distributor.

The coefficient category $\W$ is complete, so its slice $\W/I$ is finitely complete. For a finite diagram in $\Fam(\W/I)$, take the finite limit $Z$ of its raw finite carriers. For each compatible tuple $z\in Z$, take the limit of the corresponding coefficient diagram in $\W/I$. A cone factors uniquely on raw carriers and then uniquely on each coefficient. This proves the finite-limit universal property and shows that the carrier functor preserves it. The terminal diagram is included: its carrier is a singleton and its coefficient is the terminal object of the slice. Equalizers may have empty carriers, with the unique augmentation from the initial additive object.

Full faithfulness and strong preservation of the original dependent/occurrence structure are those of the conservative coefficient envelope. They do not require the new finite limits to have belonged to the original source.
\end{proof}
The finite-limit theorem concerns unrestricted augmented families. At an external input interface $A$, the unrestricted finite-fibre context category is pointwise $\Aug(\W)^A$ (for example, take the original admissibility class $\mathcal J$ to be all set maps). If a smaller original class $\mathcal J$ is retained, closure of the constructed limit projection in $\mathcal J$ is an additional requirement; it does not follow from pullback and composition stability alone. Within the unrestricted context category, reindexing along a pure input map preserves finite limits. The extension carrier over $A$ therefore preserves these limits in $\mathbf{Set}/A$. The existing local joint-comparison manuscript already constructs joint grades from coefficient limits and uses the augmented-vector-space specialization of this finite-family slice description. The increment here is a canonical conservative enlargement supplying the context-side limits for the stated entire source class, not a newly invented finite-limit constructor. Its writer-family content hypotheses are separate: compatible content-side limits and extension preservation must still be supplied, and the writer theorem does not automatically cover arbitrary original polynomial content. To combine a newly constructed limit with the original crossing, its projection must remain in the admitted class, or the content cartesian/display hypotheses must independently be extended to it. Widening the context class alone does not extend a contentad supplied only over the old class. No preservation of all joint limits by dependent accumulation or tensor is inferred.

\section{Why coproductification cannot be imposed on every old rig}
\begin{lemma}[Full-faithfulness obstruction]
If a fully faithful functor $F:\V\to\mathcal E$ is strong for the specified additive monoidal structures and the target additive tensor is categorical coproduct, then the original specified additive tensor is itself categorical coproduct. Its additive unit is initial.
\end{lemma}
\begin{proof}
For old objects $A,B,C$, strongness and full faithfulness give
\[
 \V(A\oplus B,C)
 \cong\mathcal E(FA\amalg FB,FC)
 \cong\V(A,C)\times\V(B,C).
\]
The target injections lift uniquely to old arrows by fullness, and faithfulness reflects their universal equations. The nullary version gives a singleton $\V(0,C)$. Thus the original sum and zero already have the asserted universal properties.
\end{proof}
For the discrete nonnegative-real rig, addition is not categorical coproduct: for example the canonical injection from coefficient $1$ to coefficient $1+1$ does not exist. Therefore the preceding positive construction cannot be forced onto this source while simultaneously retaining full faithfulness and strong preservation of its given additive tensor. This does not classify every possible joint completion; it excludes one tempting uniform repair with its exact violated preservation obligation.

\section{Position relative to the full target}
The positive theorem constructs an actual native context-limit source for all original rigs satisfying the stated additive condition, rather than merely assuming coefficient limits will suffice. The selected-law theorem excludes a weak product throughout the different scalar envelope. These statements explain why the same general coefficient enlargement can preserve O1 interaction yet fail to supply a universal shared probabilistic context.

Neither conclusion says that shared randomness is impossible. A specified joint context, a category of alternative couplings, or a law specification can retain it. A general completion must decide which of those interfaces is intended and prove compatibility with native substitution and feedback. The original fully faithful embeddings, directed native cells, and finite scope data must not be silently replaced by output-law equivalence. General feedback, arbitrary joint-scope substitution, and infinite-effect completion remain open.

\begin{thebibliography}{9}
\bibitem{Envelope} \emph{Directed distributivity and a conservative native coefficient envelope}, accompanying project manuscript. Original native-map preservation and opposite Day construction.
\bibitem{O1} \emph{General Independent Interaction Criterion}, project source, 1 November 2026. The native aggregation/augmentation presentation.
\bibitem{GLF} R. Garner and I. L\'opez Franco, \emph{Commutativity}, Definition 42 and Proposition 43. \url{https://arxiv.org/pdf/1507.08710}.
\bibitem{Joint} \emph{Local joint comparison} and \emph{Local product shared interaction}, project manuscripts in \path{foundations/local-joint-comparison/} and \path{foundations/shared-witness-interaction/}. Existing native context/content limit and projection-derived interaction machinery; its content and preservation hypotheses remain distinct from the context-side result here.
\end{thebibliography}
\end{document}
